\section{Related Work}
There is a long history of pre-training general language representations, and we briefly review the most popular approaches in this section.
\subsection{Feature-based Approaches}
Learning widely applicable representations of words has been an active area of research for decades, including non-neural~\cite{brown-etal:1992:_class,ando-zhang:2005, blitzer-mcdonald-pereira:2006:_domain} and neural~\cite{collobert-weston:2008, mikolov-etal:2013, pennington-socher-manning:2014:_glove} methods. Pre-trained word embeddings are considered to be an integral part of modern NLP systems, offering significant improvements over embeddings learned from scratch~\cite{turian-ratinov-bengio:2010:_word_repres}.

These approaches have been generalized to coarser granularities, such as sentence embeddings~\cite{kiros-etal:2015:_skip, logeswaran2018an} or paragraph embeddings~\cite{le-mikolov:2014:_distr}. As with traditional word embeddings, these learned representations are also typically used as features in a downstream model.

ELMo~\cite{peters-etal:2017:_semi} generalizes traditional word embedding research along a different dimension. They propose to extract \emph{context-sensitive} features from a language model. When integrating contextual word embeddings with existing task-specific architectures, ELMo advances the state-of-the-art for several major NLP benchmarks~\cite{peters-etal:2018:_deep} including question answering~\cite{rajpurkar-etal:2016:_squad} on SQuAD, sentiment analysis~\cite{socher-etal:2013:_recur}, and named entity recognition~\cite{tjong-de:2003}.

\subsection{Fine-tuning Approaches}
A recent trend in transfer learning from language models (LMs) is to pre-train some model architecture on a LM objective before fine-tuning that same model for a supervised downstream task~\cite{dai-le:2015:_semi, howard-ruder:2018, radford-etal:2018}. The advantage of these approaches is that few parameters need to be learned from scratch. At least partly due this advantage, OpenAI GPT~\cite{radford-etal:2018} achieved previously state-of-the-art results on many sentence-level tasks from the GLUE benchmark~\cite{wang-etal:2018:_glue}.

\begin{figure*}[t]
\begin{center}
\includegraphics[width=\textwidth]{figures/BERT_comparisons.pdf}
\end{center}
\caption{Differences in pre-training model architectures. BERT uses a bidirectional Transformer. OpenAI GPT uses a left-to-right Transformer. ELMo uses the concatenation of independently trained left-to-right and right-to-left LSTM to generate features for downstream tasks. Among three, only \bert representations are jointly conditioned on both left and right context in all layers.}
\label{fig:BERT_comparisons}
\end{figure*}

\subsection{Transfer Learning from Supervised Data}
While the advantage of unsupervised pre-training is that there is a nearly unlimited amount of data available, there has also been work showing effective transfer from supervised tasks with large datasets, such as natural language inference~\cite{conneau-EtAl:2017:EMNLP2017} and machine translation~\cite{mccann-etal:2017:_learn_trans}. Outside of NLP, computer vision research has also demonstrated the importance of transfer learning from large pre-trained models, where an effective recipe is to fine-tune models pre-trained on ImageNet~\cite{imagenet_cvpr09, yosinski2014transferable}. 
